\section{Brother, Pastor, and Defender of Nicene Faith}

Ambrose's first decade as bishop (374--383) was forged in a crucible of domestic bereavement, ascetical controversy, barbarian catastrophe, and Trinitarian conflict. Readers who know Ambrose only as the stern churchman who confronted emperors often miss the tenderness that fills his earliest writings---above all, the two orations he delivered in Milan on the death of his brother Satyrus.

\subsection{The Death of Satyrus and the Hope of Resurrection}

Satyrus died suddenly after returning to Italy from an arduous journey to Sicily and Africa to recover family debts wrongfully withheld by a certain Prosper (\work{De Excessu Fratris Satyri} I.17, 24, 26). On the sea voyage he survived a shipwreck in an old and leaky vessel. Though not yet baptized---still ``before being initiated in the more perfect mysteries''---Satyrus feared departing life without the Eucharist, asked baptized fellow passengers for the consecrated Sacrament, bound it in a napkin around his neck, and cast himself into the waves without seeking a wooden plank (\work{De Excessu} I.43). Safely ashore, he sought baptism at once, yet refused to receive it until he had questioned the local bishop ``whether he agreed with the Catholic bishops, that is, with the Roman Church'' (\work{De Excessu} I.47). Finding the local see caught in the Luciferian schism, Satyrus waited until he could reach Catholic communion rather than accept baptism in a divided body.

Not long after reaching Milan, Satyrus succumbed to illness---dated by de Romestin to 17~October 379, while the \work{Martyrologium Romanum} keeps his deposition on 17~September and modern scholars often place the death in 378 amid the Gothic crisis after Adrianople. Ambrose preached the first oration beside the body and the second, on the Resurrection, seven days later. Nothing in Latin patristic literature surpasses the candor of his grief. Ambrose recalled how the two brothers looked so alike that citizens constantly mistook one for the other (\work{De Excessu} I.38), how they shared an undivided patrimony because Satyrus insisted ``that love was no reason for making a will'' (I.59), and how Ambrose tried to draw in his dying brother's last breath:

\begin{quote}
It profited me nothing to receive thy last breath, nor to have breathed into the mouth of thee dying, for I thought that either I myself should receive thy death, or should transfer my life to thee. O that sad, yet sweet pledge of the last kiss! O the misery of that embrace, in which the lifeless body began to stiffen, the last breath vanished! (\work{De Excessu Fratris Satyri} I.19)
\end{quote}

Yet he refused to let grief become despair. ``Tears, then, are marks of devotion, not producers of grief. I confess, then, that I too wept, but the Lord also wept. He wept for one not related to Him, I for my brother'' (I.10). Seven days later, in Book~II, Ambrose turned his own sorrow into catechesis on the Resurrection: ``What grief is there which the grace of the Resurrection does not console? What sorrow is not excluded by the belief that nothing perishes in death?'' (II.3). He buried Satyrus at the left hand of the martyr St.~Victor in Milan, composing the four-line elegiac epitaph reported in the tradition that the holy martyr's blood might bathe his brother's resting place.

\subsection{Consecrated Virginity and Civic Backlash}

Around 377, even before Satyrus's death, Ambrose published three books \work{Concerning Virgins} (\work{De Virginibus}) addressed to Marcellina, followed soon afterward by \work{Concerning Widows} (\work{De Viduis}) and \work{De Virginitate}. Opening \work{De Virginibus} on the feast of the twelve-year-old Roman martyr St.~Agnes (21~January), Ambrose presented consecrated virginity not as a Stoic discipline or a repudiation of marriage, but as a fruit of the Incarnation: ``Then a Virgin conceived, and the Word became flesh that flesh might become God'' (\work{De Virginibus} I.3.11).

His preaching provoked fierce domestic friction among noble families in northern Italy, where aristocratic marriage alliances secured estates and lineages. Ambrose himself recorded the complaints in \work{De Virginibus} I.11: critics protested that he was ``always singing the praises of virgins,'' while young women traveled from Placentia, Bononia (Bologna), and even Mauritania to receive the veil at Milan. Mothers---``and, which is more serious, mothers who were widows''---tried to bar their daughters from listening to him, drawing Ambrose's pointed challenge: ``For if your daughters desired to love a man, they could, by law, choose whom they would. Are they, then, who are allowed to choose a man not allowed to choose God?'' (\work{De Virginibus} I.11.58). He even recounted a noble maiden who fled her relatives' marriage pressures by placing her head beneath the altar itself, declaring that no bridal veil could match the altar on which Christ is consecrated (I.12.65). In \work{De Viduis}, Ambrose balanced his praise of virginity with an affirmation of honorable Christian widowhood and second marriage as lawful, setting himself apart from Montanist or Novatianist rigorism.

\subsection{Adrianople and the Living Vessels of Christ}

On 9~August 378, the Gothic victory at Adrianople annihilated the Eastern field army and killed the Homoian Emperor Valens. Refugees and captives flooded across Illyricum toward northern Italy; Ambrose's funeral oration for Satyrus shudders at the thought that ``our safety wholly depended on the barrier of the Alps'' (\work{De Excessu} I.31). When funds ran dry to ransom Christian captives from barbarian slavery, Ambrose broke up and melted down the Church's sacred gold and silver altar vessels to buy their freedom.

Homoian enemies promptly accused him of sacrilege. Years later, in \work{De Officiis} II.28, Ambrose defended his action with a theology of Christian charity that became a permanent landmark of Catholic social teaching:

\begin{quote}
Who can be so hard, cruel, iron-hearted, as to be displeased because a man is redeemed from death, or a woman from barbarian impurities, things that are worse than death, or boys and girls and infants from the pollution of idols, whereby through fear of death they were defiled? \ldots\ The Church has gold, not to store up, but to lay out, and to spend on those who need. \ldots\ It had been better to preserve living vessels than gold ones. \ldots\ The sacraments need not gold, nor are they proper to gold only---for they are not bought with gold. The glory of the sacraments is the redemption of captives. Truly they are precious vessels, for they redeem men from death. That, indeed, is the true treasure of the Lord which effects what His blood effected. (\work{De Officiis} II.28.136--138)
\end{quote}

Recalling the Roman deacon St.~Lawrence pointing to the poor and declaring, ``These are the treasures of the Church'' (II.28.140), Ambrose set careful rules for such emergencies: uninitiated vessels should be used first, broken up rather than sold intact to profane hands, and used only for urgent relief of captives, the poor, and the dead.

\subsection{Gratian, \work{De Fide}, \work{De Spiritu Sancto}, and Aquileia (378--381)}

On the eve of marching east to aid Valens in 378, the nineteen-year-old Western Emperor Gratian asked Ambrose for a short treatise proving the full divinity of the Son against Homoian Arianism. Ambrose answered with the first two books of \work{De Fide} (\work{Exposition of the Christian Faith}), opening with a brisk pledge: ``I will take in hand a bold enterprise, yet modestly withal, not so much reasoning and disputing concerning the Faith as gathering together a multitude of witness'' (\work{De Fide} I.prol.4). After Gratian returned from Thrace and requested a fuller treatment, Ambrose added Books~III--V (c.~379--380), drawing heavily on his own sermons to the Milanese congregation to answer Homoian scriptural objections one by one. In Book~IV, while explaining John~6 against those who reduced Christ's divine nature by pointing to his human flesh, Ambrose articulated an unambiguous confession of Eucharistic transformation: ``Now we, as often as we receive the Sacramental Elements, which by the mysterious efficacy of holy prayer are transformed into the Flesh and the Blood, `do show the Lord's Death'\,'' (\work{De Fide} IV.10.125).

In 381, again at Gratian's request, Ambrose completed three books \work{On the Holy Spirit} (\work{De Spiritu Sancto}), arguing that the Holy Spirit shares one undivided substance, power, and worship with the Father and the Son. Ambrose drew openly on Greek theological sources---especially Didymus of Alexandria's treatise on the Spirit, St.~Basil, and St.~Athanasius---adapting their arguments into fluent Latin for Western readers who lacked Greek. That heavy reliance provoked a sharp rebuke from St.~Jerome---who had translated Didymus himself and charged that Ambrose had turned ``good Greek into bad Latin'' (\latin{ex Graecis bonis Latina non bona}, reported through H.~B. Swete in Blomfield Jackson's NPNF~II.8 introduction and Rufinus's \work{Apologia} cited in de Romestin's Prolegomena)---even though Jerome elsewhere saluted Ambrose as ``a pillar of the Church.'' Augustine judged Ambrose's accomplishment differently: in \work{On Christian Doctrine} IV.21.46, he quoted the opening of \work{De Spiritu Sancto}---Ambrose's typological exposition of Gideon's rock and fleece---as a model of the plain, proof-driven style of Christian teaching.

Ambrose paired theological argument with decisive conciliar action. Around 380, despite fierce opposition from the Homoian Empress-mother Justina, he intervened in Illyricum to secure the consecration of the Nicene bishop Anemius at Sirmium. In September 381, while the Eastern bishops gathered at the First Council of Constantinople, Ambrose steered the Western Council of Aquileia. Two veteran Illyrian Homoian bishops, Palladius of Ratiaria and Secundianus of Singidunum, had hoped for a broad general council where Eastern sympathizers might back them; Ambrose persuaded Gratian to restrict attendance to neighboring bishops willing to come. At Aquileia, Ambrose led the cross-examination of Palladius and Secundianus directly from the floor, securing their condemnation and deposition as recorded in the synodal letters (\work{Epistles} 9--12, summarized in de Romestin's Prolegomena). By the end of 381, Homoian Arianism had lost its intellectual and institutional foothold in northern Italy, surviving chiefly inside the imperial court itself.
